\begin{table}[H]
\caption{USGS Production Availability Measures}
\label{tab:production-availability}
\centering
\begingroup
\footnotesize
\setlength{\tabcolsep}{4pt}
\renewcommand{\arraystretch}{1.05}
\begin{tabularx}{\textwidth}{>{\raggedright\arraybackslash}p{0.09\textwidth}>{\raggedright\arraybackslash}p{0.23\textwidth}>{\raggedleft\arraybackslash}p{0.14\textwidth}>{\raggedleft\arraybackslash}p{0.14\textwidth}>{\raggedright\arraybackslash}X}
\toprule
HS code & Product & 2023 \(A^{prod}\) & 2024 \(A^{prod}\) & Mapping caveat \\
\midrule
2504 & Natural graphite & 0.216 & 0.219 & Direct production proxy \\
8110 & Antimony metal & 0.519 & 0.494 & Upstream production proxy \\
8106 & Bismuth metal & 0.152 & 0.166 & Refinery and byproduct caveat \\
284180 & Tungstates/APT & 0.176 & 0.182 & Upstream tungsten proxy \\
\bottomrule
\end{tabularx}
\caption*{\footnotesize Notes: \(A^{prod}=Scale^{prod}\times(1-HHI^{ROW,prod})\). Detailed production, scale, HHI, diversity, and country rows are reserved for the thesis appendix; the main proposal uses this compact version to avoid overcrowding. Source: author calculations from USGS \emph{Mineral Commodity Summaries 2025} world production data \citep{usgs_mcs_2025}.}
\endgroup
\end{table}
